\section{The number of ordered factorizations: incidence without replacement}
\label{sec:incidence}
For this section write $U_d$ for the uniform law on the $N_d$ simple
$d$-regular boards on the two prescribed labelled shores, $W(A)=\operatorname{per}A$,
$m_d=\mathbb E_{U_d}W$, and $\widehat U_d(A)=W(A)U_d(A)/m_d$ for the
perfect-matching-biased law. Write $t_d=\|U_d-\widehat U_d\|_{\rm TV}$
and $V_d=n!(d/n)^n$. These probability spaces are used only to count:
the sampling algorithm never needs a uniform-board or a permanent-counting oracle.
We shall prove a constant-factor lower bound, rather than a power-of-$n$
lower bound which would change the running-time exponent.

\begin{proposition}[Constant-loss count lower]\label{prop:count-lower}
For an absolute $c_*>0$ and all sufficiently large $n$,
\[
 |\mathcal L_n|\ge c_*\prod_{d=1}^n m_d .
\]
\end{proposition}

\paragraph{Killing a marked incidence.}
Counting incidences $(A,M)$ with $A$ a $d$-board and $M$ one of its
matchings gives $N_d m_d=N_{d-1}m_{n-d+1}$, by deleting $M$.
At $A$ choose a uniform matching $M$, put $T=A-M$, and retain the step
with probability
\[
 \min\{1,W(A)/m_d\}\min\{1,m_{n-d+1}/W(K_{n,n}-T)\}.
\]
Starting from any subprobability dominated by $U_d$, the first factor
discards mass at most $t_d$, and the surviving marked incidence is
dominated pointwise by the uniform incidence. Its $T$ marginal is
therefore bounded by $U_{d-1}(T)W(K_{n,n}-T)/m_{n-d+1}$.
The second factor discards at most $t_{n-d+1}$ and restores domination
by $U_{d-1}$. Each particular matching transition now has probability
at most $1/m_d$. We need the following three estimates:
for fixed large $c$, with $D_0=\lfloor n^{1/10}\rfloor$ and
$R=\lceil\sqrt n\,\ell^{B_0}\rceil$ for a sufficiently large fixed $B_0$,
\[
 \sum_{c<d\le D_0}t_d\le C/c+o(1),\qquad
 \sum_{D_0<d<R}t_d=o(1),\qquad \sum_{R\le d\le n}t_d=o(1).
\]
The cutoff $D_0$ is internal to this incidence proof.
\paragraph{Sparse strip.}
For $c\le d\le D_0$, enumeration of simple bipartite boards with
prescribed sparse margins and forbidden cells gives
\begin{align}
 m_d&=n!d^{2n}\frac{((d-1)n)!}{(nd)!}
                      \exp[-1/2+O(d^3/n)],\notag\\[-1mm]
 t_d&\le C\{d^{-2}+d^{3/2}/\sqrt n+(nd)^{-1/4}\}.\label{eq:inc-sparse}
\end{align}
For details, in the formulas of Greenhill et al.\ \cite[Theorems~2.1--2.2]{GHIM}, let
$k_i,l_j\in\{d,d-1,d-2\}$, $s=\sum k_i=\sum l_j=\Theta(nd)$,
$S_2=\sum(k_i)_2$, $T_2=\sum(l_j)_2$, and
$F_X=\sum_{ij\in X} k_i l_j$ for a forbidden board $X$ of maximum
degree two. Their enumeration in this range reduces uniformly to
\[
 \frac{s!}{\prod k_i!\prod l_j!}
 \exp[-S_2T_2/(2s^2)-F_X/s+O(d^3/n)].
\]
Indeed $d^2=o(s^{2/3})$, their forbidden parameter is $O(d^2)=o(s)$,
$F_X=O(nd^2)=o(s^{5/3})$, and $O(d^2)F_X=o(s^2)$.
For one compulsory matching apply the formula to its complement after
decrementing its endpoint degrees by one. For two specified matchings
with $j$ common edges set $s=n(d-2)+j$ and $u=(d-1)j/s$.
Directly $S_2/s=T_2/s=d-3+u$ and
$F_X/s=2(d-2)-(d-3)u$, giving a net exponential relative to the
unrestricted count of $\exp(-u^2/2+O(d^3/n))$. There are
$(n!)^2 h_{n-j}/j!$ pairs where $h_m$ is the fraction of derangements
of $m$ objects. Set $x=(d-1)^{-2}$, $s_0=n(d-2)$ and let $K$ have
mass $\binom{s_0+j}{j}(1-x)^{s_0+1}x^j$. Extracting the factorials
and using Stirling yields
\[
 \mathbb E_{U_d}(\operatorname{per}A)^2/m_d^2
 =(1-x)^{-1/2}(1+O(d^3/n))
 \mathbb E[\mathbf1_{K\le n}e h_{n-K}e^{-u(K)^2/2}].
\]
Here $\mathbb EK=n/d+O(1)$, $\operatorname{Var}K=O(n/d)$, $K>n/2$
has exponentially small probability, $u(n/d)=1/(d-1)$ and
$|u'|=O(1/n)$. The expectation is $e^{-x/2}+O(1/\sqrt{nd})$.
Since $(1-x)^{-1/2}e^{-x/2}=1+O(x^2)$, total variation is bounded
by half the coefficient of variation, as required. Stirling in the
first line gives
$\log(m_d/(L_d\sqrt{n/d}))=O(d^{-1}+d^3/n)$ also for $2\le d<c$:
only the one-mark enumeration is used there. At $d=1$, $m_1=L_1=1$.

\paragraph{Moderate strip.}
Take $D_0<d<R=\lceil\sqrt n\,\ell^{B_0}\rceil$, $B_0$ fixed sufficiently
large, and $p=d/n$, $\tau=1/(nd^{0.9})+d/n^2+p^3$.
The bipartite case of Liebenau--Wormald \cite[Theorem~1.4]{LW} implies that for any fixed embedded simple pattern $F$
with $h$ edges, degrees $h_v$, and $w(F)=\sum_v\binom{h_v}{2}$,
\[
 \Pr_{U_d}(F\subset A)
 =\frac{\prod_v(d)_{h_v}}{(nd)^h}\exp\{w(F)/n+O_F(\tau)\}.
 \tag{I1}\label{eq:inc-pins}
\]
This is not a direct invocation with forbidden cells. After any bounded
number of margin decrements the unrestricted edge marginal is
$q_{ij}=(s_i t_j/m)(1+O_F(1/(nd^{0.9})))$: in that theorem take
$\kappa=.01$, $m=nd-O(1)$, endpoint deviations $O(1)$ and margin
variance $O(1/n)$, at density $o(1)$. Decrement the edges successively
using the exact ratio of margin counts $q_e^{\rm before}/(1-q_e^{\rm after})$,
then multiply by the probability of avoiding $F$ in the final reduced
law. In that law any two fixed pins have probability
$p^2(1+O_F(1/d+p))$ and three cost $O_F(p^3)$ by the same ratios
and union bounds. Thus log avoidance is
$-\sum_{e\in F}q_e^{\rm final}-hp^2/2+O_F(p^3+p^2/d)$.
The inverse denominators cancel the quadratic term.
In $\sum(q_e^{\rm after}-q_e^{\rm final})$ each later shared endpoint
contributes $1/n$, giving $w(F)/n+O_F(\tau)$. In particular
disjoint patterns factor to relative $1+O_F(\tau)$.

For connected unicyclic $F$ let $Y_F=\operatorname{inj}(F,A)/d^h$.
Its mean is $1-w(F)/d+O_F(d^{-2}+n^{-1}+\tau)$ and its variance is
$O_F(\tau+(nd)^{-1}+d^{-4})$. Indeed disjoint placements use \eqref{eq:inc-pins};
shared vertices but no edges change the normalized weight by
$O_F(1/d)$ on an $O_F(1/n)$ fraction (merge the vertex falling
factorials). For $j\ge1$ shared edges on $u$ shared vertices the
joint bound is $O_F(n^{j-u}/d^j)$; either $u\ge j+1$, or the
common graph contains the cycle, hence $j\ge4$ and $u\ge j$.
The absolute mean deviation from $m_F=1-w(F)/d$ is summable $o(1)$
over this strip, since in particular $\sum\sqrt\tau=o(1)$.

For $\widehat U_d$, delete the distinguished matching to obtain a $(d-1)$-board
tilted by its number of \emph{complement} matchings, then add its
conditional uniform complement matching. The tilt density is at most
$2$ by Br\'egman/van der Waerden since $p\log n=o(1)$.
For fixed $j$ disjoint available pins the complement matching probability
is $n^{-j}(1+O_j(p+1/n))$: switching two matched rows gives $n-O(d)$
available switches out of a pin; the unique reverse insertion fails
only if a forbidden cross-edge is present, with probability $O(d/n)$
by the two-pin upper $O(n^{-2})$ (the same upper applies after deleting
a matched row and column). The part of $Y_F$ with no new edge is thus
within summable absolute-mean error of $1-[w(F)+h]/d$.
For one specified edge marked, let $Y_e$ be its contribution. Under
the untilted remainder the first two moments of $dY_e$ give
$\mathbb E(dY_e-1)^2=O_F(1/d+p)$. To see this, disjoint embeddings in
the second moment contribute at most $1+O_F(1/d+p)$ by \eqref{eq:inc-pins} and the
two-mark bound. An overlap on $u$ vertices, $j$ old edges and $k=0,1$
common marks costs $O_F(n^{j+k-u}/d^j)$: if cyclic, $j\ge3$, otherwise
$u\ge j+k+1$. The first moment is $1+O_F(1/d+p)$ by the one-mark
estimate and a conditional absence probability $1-O_F(p)$ on the marked
cell. The nonnegative squared-error bound survives the bounded tilt.
Hence each singleton mark adds $1/d$ with summable absolute-mean error
$O_F(d^{-1}\sqrt{1/d+p})$. Terms with two or more marks cost
$O_F(d^{-2})$. This proves the same $m_F$ centering under both laws.

Apply the finite connected-template expansion of the log rook polynomial
before its Gamma averaging, at a fixed order $s>30$ (enlarge the order
if necessary so $\sum n\ell^{O(s)}/d^s=o(1)$). In each template delete
the $j$ constant $J$ edges; if $h'$ components remain, $j\ge h'-1$.
Expand homomorphisms into collision partitions, retaining tree
homomorphism terms exactly ($nd^e$), and cyclic simple injections.
A single unicyclic injection is replaced by $m_F$ after normalizing
by its number of distinct edges; retain the original further powers
of $1/d$ for repeated edges. Rank at least two costs $O_F(1/n)$ in
absolute mean by \eqref{eq:inc-pins} or the planted-mark decomposition. With $w\ge2$
cyclic components the remaining tree factors supply at most $n^{1-w}$;
two normalized spanning unicyclic counts have bounded joint mean
by the overlap argument above, and each further one is at most $n/d$.
Thus these terms cost $O_F(1/n)$. Only finitely many templates are
involved. Let $P_0$ be the resulting deterministic polynomial and set
\[
 c_d=\log\mathbb E[\mathbf1_{X\in I}\exp P_0(1/X)],\qquad
 X\sim\operatorname{Gamma}(n+1,n),\quad I=[1-1/(2\ell),1+1/(2\ell)]
\]
(shape, rate). On healthy boards the finite polynomial $P_A$ approximates
the actual log rook polynomial on the inverse window; the discarded
signed Gamma integral is $e^{-\Omega(n/\ell^2)}$ in absolute value,
negligible in log since $\operatorname{per}A/V_d\ge1$.
Theorem~B of Tikhomirov and Youssef \cite{TY} for the directed
loops-allowed model, gives the required centered singular bound with
arbitrarily high polynomial confidence under $U_d$; the planted
remainder density bound transfers it to $\widehat U_d$. Bad events have
polynomial log cost. Sandwiching the two positive integrals by
$\exp(\pm\sup_I|P_A(1/x)-P_0(1/x)|)$ gives
\[
 \sum_{D_0<d<R}(\mathbb E_{U_d}+\mathbb E_{\widehat U_d})
     |\log(\operatorname{per}A/V_d)-c_d|=o(1).\tag{I2}\label{eq:inc-common}
\]
This log argument does not require positive exponential moment control
of the error on the exceptional boards.

\paragraph{Dense strip.}
For $R\le d\le n$ let $z=d(d-1)/(n-1)$,
$c_{ij}=|N(i)\cap N(j)|$,
$W=\sum_{i<j}(c_{ij}-z)^2$, and put
\[
 w_*=(n/2)(d-z)(d-1)(1-(d-1)/(n-2)).
\]
Under $U_d$ or $U_d$ conditional on a fixed planted matching, resplit
the exclusive columns of a uniformly chosen row pair into equal halves
(subject to the at most two exclusive pins). This gives an exchangeable
pair $W,W'$. Let $2t$ be the exclusive-set size, $b_k$ its overlap
with row $k$, $D_k=c_{ik}-c_{jk}$. Then
$W'-W=\sum_{k\ne i,j}((D_k')^2-D_k^2)/2$; the subtracted half
averaged over pairs is $2W/n$. Unpinned,
$\mathbb E(D'_k)^2=b_k(2t-b_k)/(2t-1)$. Its sum differs from
$2t(d-1)(1-(d-1)/(n-2))$ by $O(d)$ on the good spectral event:
the squared deviations of the $b_k$ from their mean total $O(dt)$
by applying the centered board to the exclusive-set indicator.
Forcing up to two columns changes the sum by $O(d)$; for $t>2$ the
hypergeometric conditional variance and squared-mean changes per row
are $O(r_k+b_k/t)$, where $r_k$ counts pins incident to row $k$.
Smaller $t$ is immediate. Thus the drift is $-2(W-w_*)/n+O(d)$.
With both endpoints spectrally healthy, $|W'-W|=O(d^2)$ by applying
the centered norm bound to the difference of the pair's row indicators.
Bad endpoints have arbitrarily small polynomial mass under either law:
for dense $d$ use complementing, and transfer the plant by bounded
polynomial Br\'egman density (for $d<n/2$ via its remainder).
The identity
$-\mathbb E[(W-w_*)(W'-W)]=\mathbb E(W'-W)^2/2$
now yields $\mathbb E|W-w_*|=O(nd+\sqrt{n}\,d^2)$ under both laws.

The quartic signed log rook expansion has only one graph-dependent
term on the Gamma window, namely $x^{-4}W/(2d^4)$. Indeed, with
$B_A=A/d-J$, the trace quartic equals
$[nd^2+2W+n(n-1)z^2]/d^4-1$. Its remainder is
$O(n\ell^2/d^3+n^{-20})$ on healthy uniform or planted boards.
Replacing $W$ by $w_*$ in the truncated positive integral gives \eqref{eq:inc-common}
also over $R\le d\le n$ with summed bound
$O(\ell^{2-2B_0}+\ell^{-B_0})+o(1)$.


\paragraph{From the three strips to total variation.}
Write $h=\log(W(A)/V_d)$ and $M_d=\log(m_d/V_d)$.
Relative entropy in both directions gives
$\mathbb E_{U_d}h\le M_d\le\mathbb E_{\widehat U_d}h$.
Since $|e^x-1|\le(1+e^x)|x|$,
$2t_d\le(\mathbb E_{U_d}+\mathbb E_{\widehat U_d})|h-M_d|$.
The common centers in \eqref{eq:inc-common} thus prove the two upper-strip
claims; \eqref{eq:inc-sparse} proves the first since
$\sum_{d\le D_0}d^{3/2}/\sqrt n=O(n^{-1/4})$ and
$\sum_{d\le D_0}(nd)^{-1/4}=O(n^{-7/40})$.
In particular, for $c$ a sufficiently large \emph{fixed} integer, the
middle experiment starting at $U_{n-c}$ and stopping at degree $c$ retains
at least $1-2\sum_{d=c+1}^{n-c}t_d\ge 3/4$ for large $n$. Both
endpoint marginals of the surviving measure are dominated by their
respective uniform laws. No independence of the endpoints is asserted.

\paragraph{Why neither endpoint loses a power of $n$.}
Let $F(G)$ be the number of ordered $c$-factorizations of a $c$-board,
let $R_c$ be the number of $c$-row Latin rectangles, and set
$\mu_c=R_c/N_c=\mathbb E_{U_c}F$. Drawing $c$ independent uniform
perfect matchings conditional on disjointness and forgetting their order
gives exactly the density $F/\mu_c$ with respect to $U_c$. For fixed
$c\ge3$, this law and $U_c$ are mutually contiguous:
the graph-restricted sum of $c$ one-factors on prescribed labeled shores has this contiguity by Greenhill, Kim and Wormald \cite[proof of Corollary~1.1, equation~(1)]{GKW}. It is the perfect-matching contiguity in their proof, not their result on Hamilton cycles, that is used here.
Consequently for any $\varepsilon>0$ some $\delta>0$ satisfies
$U_c(F<\delta\mu_c)<\varepsilon$ eventually. Otherwise choose a
subsequence with $U_c(F<\mu_c/j)\ge\varepsilon$. The tilted masses
on it are at most $1/j$; extend the events by the empty set off the
subsequence to contradict contiguity.

Restrict the surviving middle measure to paths whose initial
\emph{complement} and final board both have $F\ge\delta\mu_c$.
With $\varepsilon=1/8$ the two removals cost at most $2\varepsilon$
even jointly, by endpoint domination. Every particular surviving
middle path (including its initial state) has probability at most
$N_c^{-1}\prod_{d=c+1}^{n-c}m_d^{-1}$, whereas every retained path
extends in at least $(\delta\mu_c)^2$ ways by ordering top and bottom
factors. A labelled complete square determines its middle path and
both endpoint orders uniquely. We obtain
\[
 |\mathcal L_n|\ge \tfrac12\delta^2\frac{R_c^2}{N_c}
                   \prod_{d=c+1}^{n-c}m_d .
\]
Finally $R_c=(1+o(1))\prod_{d=n-c+1}^n m_d$: at each of the first
$c$ steps the ratio of the Br\'egman and van der Waerden bounds on the
remaining $(n-O(c))$-board is $\exp(O_c(\log n/n))$ by Stirling,
so every prefix has the same extension factor to this accuracy.
The incidence identities telescope exactly to
$\prod_{d=n-c+1}^n m_d=N_c\prod_{d=1}^c m_d$.
Substitution proves Proposition~\ref{prop:count-lower}. The constants
in contiguity need not define a computable membership test: the two
endpoint cuts appear only in this counting argument.
